
\documentclass{article}
\usepackage{colortbl}
\usepackage{makecell}
\usepackage{multirow}
\usepackage{supertabular}

\begin{document}

\newcounter{utterance}

\centering \large Interaction Transcript for game `cladder', experiment `full\_v1.5\_default', episode 6572 with selfrefine{-}llama{-}3.1{-}8b{-}instruct.
\vspace{24pt}

{ \footnotesize  \setcounter{utterance}{1}
\setlength{\tabcolsep}{0pt}
\begin{supertabular}{c@{$\;$}|p{.15\linewidth}@{}p{.15\linewidth}p{.15\linewidth}p{.15\linewidth}p{.15\linewidth}p{.15\linewidth}}
    \# & \multicolumn{2}{c}{Player} && \multicolumn{2}{c}{Game Master} \\
    \hline

    \theutterance \stepcounter{utterance}  
    & & & \multicolumn{4}{p{0.6\linewidth}}{
        \cellcolor[rgb]{0.9,0.9,0.9}{
            \makecell[{{p{\linewidth}}}]{
                \texttt{\tiny{[P1$\langle$GM]}}
                \texttt{You are an expert in causal inference. The following question is not a typical commonsense query, but rather a meticulously designed question created by a professor specializing in causal inference, intended to assess the students' mastery of the course content.} \\
\texttt{Imagine a self{-}contained, hypothetical world with only the following conditions, and without any unmentioned factors or causal relationships: Kwox has a direct effect on kwoz. Swoq has a direct effect on kwoz. For those who are not kwox, the probability of kwoz is 82\%. For those who are kwox, the probability of kwoz is 32\%. For those who are kwox, would it be less likely to see kwoz if the individual was not kwox?} \\
\texttt{Start your answer with 'Yes' or 'No', followed by additional reasoning or evidence to support your explanation.} \\
            }
        }
    }
     \\ \\

    \theutterance \stepcounter{utterance}  
    & & \multicolumn{4}{p{0.6\linewidth}}{
        \cellcolor[rgb]{0.95,0.95,0.95}{
            \makecell[{{p{\linewidth}}}]{
                \texttt{\tiny{[P1$\leftrightarrow$P]}}
                \texttt{{[}DRAFT{]}} \\
\texttt{No.} \\
\\ 
\texttt{To answer this question, we need to consider the causal relationships and the conditional probabilities given. We are told that Kwox has a direct effect on kwoz, and Swoq has a direct effect on kwoz. However, the question does not provide any information about the relationship between Kwox and Swoq, or whether they are confounding variables.} \\
\\ 
\texttt{Given the conditional probabilities, we can see that the probability of kwoz is lower for those who are Kwox (32\%) compared to those who are not Kwox (82\%). This suggests that being Kwox has a protective effect on the occurrence of kwoz.} \\
\\ 
\texttt{However, the question asks whether it would be less likely to see kwoz if the individual was not Kwox, given that they are Kwox. This is a counterfactual question, which requires us to consider the potential outcomes if the individual were not Kwox.} \\
\\ 
\texttt{To answer this question, we need to consider the causal effect of Kwox on kwoz. If Kwox has a direct effect on kwoz, then being Kwox would reduce the probability of kwoz. However, the question is asking whether it would be less likely to see kwoz if the individual were not Kwox, given that they are Kwox.} \\
\\ 
\texttt{In this case, we can use the concept of causal mediation to answer the question. If Kwox has a direct effect on kwoz, then the causal effect of Kwox on kwoz is not fully explained by the conditional probability of kwoz given Kwox. This suggests that there may be other factors at play that are influencing the relationship between Kwox and kwoz.} \\
\\ 
\texttt{However, without further information about the relationship between Kwox and Swoq, or whether they are confounding variables, we cannot determine the causal effect of Kwox on kwoz. Therefore, we cannot conclude that it would be less likely to see kwoz if the individual were not Kwox, given that they are Kwox.} \\
\\ 
\texttt{However, we can use the concept of counterfactuals to answer the question. The counterfactual question is asking whether it would be less likely to see kwoz if the individual were not Kwox, given that they are Kwox. This is equivalent to asking whether the probability of kwoz would be higher if the individual were not Kwox, given that they are Kwox.} \\
\\ 
\texttt{Using the law of total probability, we can calculate the probability of kwoz given that the individual is Kwox and not Kwox. Let P(Kwox) be the probability of being Kwox, and P(kwoz|Kwox) be the probability of kwoz given that the individual is Kwox. Then, the probability of kwox and kwoz is} \\
            }
        }
    }
    & \\ \\

    \theutterance \stepcounter{utterance}  
    & & \multicolumn{4}{p{0.6\linewidth}}{
        \cellcolor[rgb]{0.95,0.95,0.95}{
            \makecell[{{p{\linewidth}}}]{
                \texttt{\tiny{[P1$\leftrightarrow$P]}}
                \texttt{{[}SELF{-}FEEDBACK{]}} \\
\texttt{FORMAT: ok} \\
\texttt{RULES: ok} \\
\texttt{STRATEGY: problem} \\
\texttt{REASON: The answer is overly complex and does not directly address the question, which asks for a simple yes or no response.} \\
            }
        }
    }
    & \\ \\

    \theutterance \stepcounter{utterance}  
    & \multicolumn{4}{p{0.6\linewidth}}{
        \cellcolor[rgb]{0.561,0.737,0.561}{
            \makecell[{{p{\linewidth}}}]{
                \texttt{\tiny{[P1$\rangle$GM]}}
                \texttt{No.} \\
\\ 
\texttt{The counterfactual question is asking whether it would be less likely to see kwoz if the individual were not Kwox, given that they are Kwox. This is equivalent to asking whether the probability of kwoz would be higher if the individual were not Kwox, given that they are Kwox. However, the probability of kwoz is lower for those who are Kwox (32\%) compared to those who are not Kwox (82\%), suggesting that being Kwox has a protective effect on the occurrence of kwoz. Therefore, it would not be less likely to see kwoz if the individual were not Kwox.} \\
            }
        }
    }
    & & \\ \\

    \theutterance \stepcounter{utterance}  
    & & & \multicolumn{2}{p{0.3\linewidth}}{
        \cellcolor[rgb]{0.9,0.9,0.9}{
            \makecell[{{p{\linewidth}}}]{
                \texttt{\tiny{[GM$|$GM]}}
                \texttt{No} \\
            }
        }
    }
    & & \\ \\

    \theutterance \stepcounter{utterance}  
    & & & \multicolumn{2}{p{0.3\linewidth}}{
        \cellcolor[rgb]{0.9,0.9,0.9}{
            \makecell[{{p{\linewidth}}}]{
                \texttt{\tiny{[GM$|$GM]}}
                \texttt{no} \\
            }
        }
    }
    & & \\ \\

    \theutterance \stepcounter{utterance}  
    & & & \multicolumn{2}{p{0.3\linewidth}}{
        \cellcolor[rgb]{0.9,0.9,0.9}{
            \makecell[{{p{\linewidth}}}]{
                \texttt{\tiny{[GM$|$GM]}}
                \texttt{game\_result = WIN} \\
            }
        }
    }
    & & \\ \\

\end{supertabular}
}

\end{document}
